\documentclass[11pt,a4paper]{article}

\usepackage[margin=22mm]{geometry}
\usepackage{amsmath,amssymb}
\usepackage{booktabs}
\usepackage{array}
\usepackage[hidelinks]{hyperref}
\usepackage{xurl}

\title{Beam-Spline Mappers in the Mok FSI Case\\
Theory, Current Implementation, Verification, and Stability Evidence}
\author{}
\date{24 July 2026}

\newcommand{\RR}{\mathbb{R}}
\newcommand{\T}{\mathsf{T}}
\newcommand{\curl}{\nabla\times}

\begin{document}
\maketitle

\section{Scope and canonical setup}

This note is the single maintained technical record for
\texttt{FSI\_mok\_Test\_Case}.  It explains the current implementation of
\texttt{BeamSplineMapper} and
\texttt{BeamSplineMapper\_WithRotationalRecovery}, the analytical benchmark
design, the conservative inverse maps, and the retained nonlinear FSI
evidence.  Historical parameter sweeps and pre-fix output are not part of the
current evidence set.

The physical model follows the official Kratos Mok validation case.  The
two-dimensional wall occupies
\[
  x\in[0.4965,0.5015]\ {\rm m},\qquad y\in[0,0.25]\ {\rm m},
\]
and its beam abstraction is placed on the centroidal line \(x=0.499\) m.  The
formal linear-thickness scale is \(s=50\).  It is not a numerical tuning
parameter: structure and fluid load are scaled together.  The beam properties
are
\[
\begin{aligned}
 A &=0.25, &
 I_{22}&=2.0833333333333336\times10^{-2},\\
 I_{33}&=5.208333333333334\times10^{-7},&
 J&=2.0833333333333334\times10^{-6}.
\end{aligned}
\]
Fluid \texttt{REACTION} is multiplied by \(+50\) before transfer, and
\texttt{swap\_sign} remains active.  The retained FSI gate uses
\[
 \Delta t=0.05,\qquad \alpha=0.03,\qquad
 N_{\rm coupling}=60.
\]

\section{Plain BeamSplineMapper}

\subsection{Forward map}

The plain mapper is a one-dimensional beam-spline construction.  A surface
node is projected onto the reference beam chain.  Axial translation and
torsion are interpolated on the local beam coordinate; transverse centerline
translations use spline systems whose derivative data are supplied by the
beam rotations.  A fixed affine augmentation \(\{1,\xi\}\) is used, with
\[
  \xi=\frac{x-X_c}{a}\in[-1,1],
\]
where \(X_c\) is the reference-chain midpoint and \(a=L/2\).  Derivative rows
are scaled by \(1/a\).  This normalization removes dependence on coordinate
translation and avoids the dimensional imbalance of an unscaled beam
coordinate.

The destination displacement has the generic form
\[
  u_d=F_{\rm plain}(q_b).
\]
Its section-offset contribution uses the recovered/interpolated section
rotation.  Because a finite rotation matrix is used in that contribution,
\(F_{\rm plain}\) is nonlinear in the rotation degrees of freedom even though
the spline solves themselves are linear.

\subsection{Conservative inverse}

The load map is defined by discrete virtual work:
\[
  \delta u_d=J_{\rm plain}(q_b)\,\delta q_b,\qquad
  Q_b=J_{\rm plain}(q_b)^\T f_d.
\]
The bending contribution therefore cannot be distributed only to the two
nodes of the projected segment.  Its evaluation functional is first
accumulated over all destination forces and is then propagated to the whole
beam support through the transpose spline solve.  The normalized derivative
scaling used in the forward map is also used in this adjoint.

\paragraph{Resolved correctness issue.}
The earlier inverse path used only a two-node force/lever-arm distribution and
was not the transpose of the global spline interpolation.  It also omitted
part of the section offset induced by bending rotations.  The repaired map
uses the exact forward tangent and its transpose.  The current four
directional virtual-work errors are
\[
  9.88\times10^{-10},\quad 4.52\times10^{-10},\quad
  1.41\times10^{-10},\quad 2.40\times10^{-9}.
\]
All pass the \(10^{-7}\) tangent/adjoint acceptance criterion.

\section{BeamSplineMapper with rotational recovery}

\subsection{Small-rotation field}

Ahrem, Beckert, and Wendland reconstruct one vector displacement field
\(s:\RR^3\rightarrow\RR^3\).  At structural sites \(X_i\), the generalized
interpolation conditions are
\[
  s(X_i)=g_i,\qquad \curl s(X_i)=2h_i.
\]
The second equation follows from infinitesimal kinematics,
\[
  h=\frac12\curl u,
\]
and the paper explicitly restricts the interpretation to small angles.  With
three value and three curl functionals per point, the saddle system is
\[
\begin{bmatrix}
 A_{\Phi,X}&P_X\\
 P_X^\T&0
\end{bmatrix}
\begin{bmatrix}a\\b\end{bmatrix}
=
\begin{bmatrix}g\\2h\\0\end{bmatrix}.
\]
In \texttt{rotation\_recovery\_mode = small}, the surface displacement is
evaluated directly:
\[
  u_s=s(X_s).
\]
For fixed support this map is linear, \(u_s=Hq_b\), so the conservative load
map is \(Q_b=H^\T f_s\).

\subsection{Polynomial levels}

The public integer setting \texttt{polynomial\_level} has the following
meaning:
\[
\begin{array}{ccl}
0&:&\text{automatic rank/conditioning fallback},\\
1&:&\text{paper full-3D affine space, 12 vector modes},\\
2&:&\text{line-basic space, 9 modes},\\
3&:&\text{line-enriched space, 12 modes},\\
4&:&\text{centered/scaled Legendre high-order space, 30 modes}.
\end{array}
\]
Level 4 replaces dimensional monomials by Legendre polynomials in
\(\xi\in[-1,1]\) without reducing the polynomial span.  Level 1 can be
non-unisolvent on a perfectly collinear value/curl support and is rejected
when its functional matrix is rank deficient.  Level 0 selects the highest
full-rank level below the conditioning limit; it selected level 4 with a
polynomial condition estimate of approximately \(1.86\times10^2\) in the
analytical beam test.

Legacy string-valued \texttt{polynomial\_basis} values are translated with a
deprecation warning.  Supplying both interfaces is an error.

\subsection{Kernel, coordinates, and regularization}

For the Kratos Gaussian convention
\[
  \phi(r)=\exp\!\left[-\frac12(r/h)^2\right],
\]
the paper radius \(d\) in \(\exp(-r^2/d^2)\) satisfies \(d=\sqrt{2}h\).
All recovery coordinates and their value/curl evaluation functionals use the
same centered reference scale.  The analytic Gaussian Hessian at \(r=0\),
\[
  \nabla\nabla\phi(0)=-\frac1{h^2}I,
\]
is retained in the curl--curl self block.  The previous generic zero-distance
branch incorrectly replaced this nonzero limit by zero.

Regularization is added only to the kernel/interpolation block.  The
polynomial constraint block remains exactly zero, as required by the saddle
formulation.  A rank defect in \(P\) is never hidden by regularizing this zero
block.  The production recovery setting is
\[
  h=0.50,\qquad \epsilon=10^{-8},\qquad
  \texttt{polynomial\_level}=4.
\]

\subsection{Finite rigid-section hybrid}

The selectable \texttt{rotation\_recovery\_mode = finite} is deliberately
labelled an Ahrem--rigid-section hybrid, not a finite-rotation theorem from
the Ahrem paper.  The Ahrem interpolation supplies a smooth displacement
field \(s:\RR^3\to\RR^3\) and an infinitesimal rotation field
\[
  \theta(X)=\frac12\curl s(X).
\]
The hybrid assumption is added afterwards: at an aerodynamic surface point
\(X_s\), the cross-section through the projected beam-center point \(X_c\)
moves as a rigid section whose finite orientation is described by the
rotation vector \(\theta_c=\theta(X_c)\).  Thus
\[
  r_0 = X_s-X_c,
  \qquad
  u_c=s(X_c),
\]
and the mapped surface displacement is
\begin{equation}
  u_s
  =
  u_c+\left[\exp([\theta_c]_\times)-I\right]r_0 .
  \label{eq:finite_hybrid_forward}
\end{equation}
Here \([a]_\times b=a\times b\).  The offset \(r_0\) is a reference offset; it
is never formed from current node coordinates.  No logarithm of a rotation
matrix is used.  This avoids the additional branch ambiguity of recovering a
rotation vector from \(SO(3)\), especially near an angle of \(\pi\).

\paragraph{Model status.}
The small-recovery mode evaluates the Ahrem field directly at the surface
node, \(u_s=s(X_s)\).  The finite hybrid instead evaluates the recovered field
at the beam-center projection \(X_c\), converts the recovered local curl into
a finite section rotation, and applies that finite section motion to the
offset \(r_0\).  Therefore the two modes are different physical models:
small mode is a general three-dimensional displacement reconstruction,
whereas finite mode imposes a rigid cross-section.  Finite mode can exactly
reproduce a rigid section rotation when the recovered \(\theta_c\) is the
correct rotation vector, but this finite interpretation is an additional beam
kinematic assumption and not a consequence of the original small-rotation
curl identity.

\subsubsection{Forward calculation workflow}

For one mapping call, the following data are available:
\[
  \{X_i,u_i,\theta_i\}_{i=1}^{N_b}
  \quad\hbox{on the beam support,}
  \qquad
  X_s
  \quad\hbox{on the surface.}
\]
The vectors \(u_i\) and \(\theta_i\) are read from the origin beam
\texttt{DISPLACEMENT} and \texttt{ROTATION} variables.  The finite hybrid is
computed as follows.

\begin{enumerate}
\item Project \(X_s\) onto the reference beam chain and obtain the projected
      center point \(X_c\).  If \(X_s\) projects onto the segment
      \((i,i+1)\) with line coordinate \(\xi\in[0,1]\), then
      \[
        X_c=(1-\xi)X_i+\xi X_{i+1},
        \qquad
        r_0=X_s-X_c .
      \]
\item Assemble the rotational-recovery Hermite interpolation system on the
      support nodes.  For each beam node there are six scalar functionals:
      three value functionals \(L_{i,a}^{u}\) and three curl functionals
      \(L_{i,a}^{\omega}\),
      \[
        L_{i,a}^{u}(s)=e_a^\T s(X_i),
        \qquad
        L_{i,a}^{\omega}(s)=e_a^\T\curl s(X_i),
        \qquad a\in\{x,y,z\}.
      \]
      The right-hand side is
      \[
        b =
        \left[
          u_1,\ldots,u_{N_b},
          2\theta_1,\ldots,2\theta_{N_b}
        \right]^\T ,
      \]
      where the factor \(2\) comes from
      \(\curl s=2\theta\).
\item Solve the saddle system
      \begin{equation}
      \begin{bmatrix}
        A & P\\
        P^\T & 0
      \end{bmatrix}
      \begin{bmatrix}
        a\\ c
      \end{bmatrix}
      =
      \begin{bmatrix}
        b\\ 0
      \end{bmatrix}.
      \label{eq:finite_hybrid_saddle_system}
      \end{equation}
      The matrix \(A\) is built by applying all value and curl functionals to
      the vector-valued kernel; \(P\) contains the selected vector polynomial
      modes.  Regularization, when used, is added only to the kernel block
      \(A\), never to the zero constraint block.
\item Evaluate the recovered center displacement and recovered center
      rotation at \(X_c\):
      \[
        u_c=s(X_c),
        \qquad
        \theta_c=\frac12\curl s(X_c).
      \]
      In matrix form this is
      \[
        u_c=E_u(X_c)
        \begin{bmatrix}a\\c\end{bmatrix},
        \qquad
        \theta_c=\frac12 E_\omega(X_c)
        \begin{bmatrix}a\\c\end{bmatrix},
      \]
      where \(E_u\) evaluates the value of the interpolant and
      \(E_\omega\) evaluates its curl.
\item Convert the rotation vector \(\theta_c\) to a rotation matrix.  Let
      \[
        \varphi=\|\theta_c\|,
        \qquad
        K=[\theta_c]_\times .
      \]
      Rodrigues' formula gives
      \begin{equation}
        R(\theta_c)
        =
        I
        +\frac{\sin\varphi}{\varphi}K
        +\frac{1-\cos\varphi}{\varphi^2}K^2 .
        \label{eq:finite_hybrid_rodrigues}
      \end{equation}
      For \(\varphi\approx0\), the code uses the corresponding Taylor
      limits, \(R\approx I+K+\frac12K^2\).
\item Apply the rigid-section offset and write the surface displacement:
      \[
        u_s=u_c+\left(R(\theta_c)-I\right)r_0 .
      \]
\end{enumerate}

This workflow is local to one destination node after the global recovery
coefficients are known.  In practice the recovery coefficients are cached per
beam chain/support and reused for all destination nodes that belong to that
support.

\subsubsection{Conservative inverse workflow}

The finite map is nonlinear because \(R(\theta_c)\) contains trigonometric
functions of the recovered rotation vector.  Its conservative inverse is
therefore the current-state tangent transpose:
\begin{equation}
  Q_b=J(q_b)^\T f_s,
  \qquad
  J(q_b)=\frac{\partial u_s}{\partial q_b}.
  \label{eq:finite_hybrid_tangent_transpose}
\end{equation}
Here \(q_b\) collects the beam nodal translations and rotations, and
\(Q_b\) contains the corresponding work-conjugate beam forces and moments.
The defining condition is discrete virtual work,
\[
  f_s^\T\delta u_s = Q_b^\T\delta q_b .
\]

For one surface force \(f_s\), variation of
\eqref{eq:finite_hybrid_forward} gives
\begin{equation}
  \delta u_s
  =
  \delta u_c
  +
  \delta R(\theta_c) r_0 .
  \label{eq:finite_hybrid_variation}
\end{equation}
The first term is transferred by the transpose of the value-evaluation
operator \(E_u(X_c)\).  The second term is transferred through the rotation
evaluation operator \(E_\omega(X_c)\) and the derivative of the exponential
map.

For a rotation vector \(\theta\), an increment \(\delta\theta\) changes the
rotated offset according to
\begin{equation}
  \delta\left(R(\theta)r_0\right)
  =
  -R(\theta)[r_0]_\times J_r(\theta)\,\delta\theta ,
  \label{eq:finite_hybrid_offset_tangent}
\end{equation}
where \(J_r\) is the right Jacobian of \(SO(3)\),
\begin{equation}
 J_r(\theta)=I-\frac{1-\cos\varphi}{\varphi^2}[\theta]_\times
 +\frac{\varphi-\sin\varphi}{\varphi^3}[\theta]_\times^2,
 \qquad
 \varphi=\|\theta\|.
 \label{eq:finite_hybrid_right_jacobian}
\end{equation}
Using \(a^\T(B\,\delta\theta)=(B^\T a)^\T\delta\theta\), the generalized
force conjugate to the recovered rotation vector is
\begin{equation}
  g_\theta
  =
  \left[-R(\theta_c)[r_0]_\times J_r(\theta_c)\right]^\T f_s
  =
  J_r(\theta_c)^\T [r_0]_\times R(\theta_c)^\T f_s .
  \label{eq:finite_hybrid_rotation_generalized_force}
\end{equation}
At zero rotation, \(J_r=I\) and \(R=I\), so
\[
  g_\theta=[r_0]_\times f_s=r_0\times f_s,
\]
which is the usual lever-arm moment.  At finite rotation, \(g_\theta\) is not
generally equal to this naive moment.  It is the component that is
work-conjugate to the selected rotation-vector parameterization.

The adjoint accumulation for the recovery coefficients can be written
compactly as
\[
  \delta W
  =
  f_s^\T E_u(X_c)\delta z
  +
  g_\theta^\T
  \left[\frac12 E_\omega(X_c)\delta z\right],
  \qquad
  z=\begin{bmatrix}a\\c\end{bmatrix}.
\]
Therefore the adjoint right-hand side at the interpolation-coefficient level
is
\begin{equation}
  r_z
  =
  E_u(X_c)^\T f_s
  +
  \frac12 E_\omega(X_c)^\T g_\theta .
  \label{eq:finite_hybrid_coefficient_adjoint_rhs}
\end{equation}
Solving the transpose of the saddle system,
\[
      \begin{bmatrix}
        A & P\\
        P^\T & 0
      \end{bmatrix}^{\T}
      y
      =
      r_z ,
\]
propagates this destination force back to the original value and curl
functionals.  The entries associated with value data become nodal force
contributions; the entries associated with curl data are multiplied by the
same factor \(2\) convention in the forward right-hand side and become
work-conjugate nodal moment contributions.  This is the algebraic reason why
the inverse map must be treated as the exact transpose of the current forward
tangent, rather than as an independent lever-arm distribution.

\subsubsection{Minimal hand calculation}

The following single-section example is sufficient to check the finite
rigid-section part by hand after \(u_c\) and \(\theta_c\) have been recovered.
Assume the beam center projection is
\[
  X_c=(0,0,0),\qquad
  X_s=(0,1,0),
  \qquad
  r_0=(0,1,0).
\]
Suppose the recovery solve gives
\[
  u_c=(0.10,0,0),
  \qquad
  \theta_c=(0,0,\alpha),
  \qquad
  \alpha=\frac{\pi}{2}.
\]
Then
\[
  R(\theta_c)=
  \begin{bmatrix}
    0&-1&0\\
    1& 0&0\\
    0& 0&1
  \end{bmatrix},
  \qquad
  R(\theta_c)r_0=(-1,0,0).
\]
The finite rigid-section displacement is
\[
  u_s
  =
  (0.10,0,0)
  +
  \left[(-1,0,0)-(0,1,0)\right]
  =
  (-0.90,-1.00,0).
\]
The corresponding small-rotation approximation would instead use
\(\theta_c\times r_0=(-\pi/2,0,0)\), giving
\[
  u_s^{small}
  =
  (0.10,0,0)+(-1.5708,0,0)
  =
  (-1.4708,0,0).
\]
This hand calculation shows the modelling difference: finite mode rotates
the reference offset as a rigid vector, whereas the small formula keeps only
the first-order term.

For the inverse of the same hand case, take a surface force
\[
  f_s=(0,1,0).
\]
At \(\alpha=\pi/2\), \(R^\T f_s=(1,0,0)\) and
\([r_0]_\times R^\T f_s=(0,0,-1)\).  Since the rotation is purely about
\(z\), the \(z\)-component of \(J_r^\T\) is unchanged, and
\[
  g_{\theta_z}=-1.
\]
Thus a positive \(y\)-force at the rotated point performs negative virtual
work against a positive additional \(z\)-rotation in this configuration.
Together with the direct translational contribution \(f_s\), this is the
quantity that must be propagated back through the transpose recovery system.
The sign follows from differentiating the finite rotated offset, not from an
ad-hoc moment rule.

\section{Common geometric and algebraic corrections}

The following defects were resolved in both mapper implementations.

\begin{enumerate}
\item Mapper search no longer writes reference coordinates into the beam
      nodes' current coordinates.  Projection may use reference geometry,
      but it has no observable side effect on the structural state.
\item When an approximate projection coordinate is clamped, its shape
      functions, projection point, distance, and section offset are rebuilt
      from that same clamped coordinate.
\item The local frame is constructed from the Cartesian axis least aligned
      with the beam tangent, followed by orthogonal projection and a
      right-handed cross product.  Negative axes, near-axis directions, and
      nearly collinear segments no longer enter discontinuous equality
      branches.
\item Reference-chain midpoint and half-length scaling are used consistently
      in values, derivatives, curls, moment functionals, and inverse
      functionals.
\item Explicit inverse matrices are not stored.  Forward systems are solved,
      and inverse functionals are aggregated per beam chain before solving
      \(A^\T z=b\).
\item \texttt{swap\_sign} is applied only at the physical load accumulation;
      it is not mixed into the mathematical transpose.
\end{enumerate}

\section{Analytical benchmark design}

Mapper accuracy is measured only against prescribed analytical fields, never
against another mapper.  For surface-node values,
\[
 e_{\rm relL2}=
 \sqrt{\frac{\sum_i\|u_i^{\rm map}-u_i^{\rm exact}\|^2}
              {\sum_i\|u_i^{\rm exact}\|^2}}.
\]
Exact-span tests with zero regularization use a practical solver-level target
of \(10^{-10}\), relaxed to \(10^{-9}\) when the assembled saddle solve and
coordinate transforms make \(10^{-10}\) marginal.  Fields that centerline
values and curl cannot uniquely identify are reported as
\emph{non-identifiable diagnostics}, not exact-reproduction failures.

\subsection{Forward verification}

The core suite contains rigid translation and infinitesimal rotation, axial
motion, bending, torsion, curl-compatible fields, polynomial-level
reproduction, nonuniform spacing, translated/scaled coordinates, negative
axes, and near-collinear geometry.  Finite-mode checks prescribe exact rigid
rotations of \(1^\circ,30^\circ,90^\circ,170^\circ\).  Every test also checks
that mapper construction leaves current coordinates unchanged.

The current rigid-rotation results are:
\[
\begin{array}{lrrrr}
\toprule
&1^\circ&30^\circ&90^\circ&170^\circ\\
\midrule
\text{BeamSplineMapper}
&5.81{\times}10^{-6}&5.22{\times}10^{-3}
&4.51{\times}10^{-2}&8.83{\times}10^{-2}\\
\text{Recovery small}
&9.18{\times}10^{-4}&2.89{\times}10^{-2}
&1.13{\times}10^{-1}&3.13{\times}10^{-1}\\
\text{Recovery finite}
&3.98{\times}10^{-15}&4.47{\times}10^{-16}
&5.71{\times}10^{-16}&1.18{\times}10^{-15}\\
\bottomrule
\end{array}
\]
These results quantify kinematic compatibility; they do not establish FSI
stability.

\subsection{Backward/adjoint verification}

For a prescribed state \(q\), perturbation \(p\), and surface force \(f\), the
centered finite-difference directional derivative is compared with the
inverse-map work:
\[
 \frac{f^\T[F(q+\varepsilon p)-F(q-\varepsilon p)]}{2\varepsilon}
 \quad\hbox{versus}\quad
 p^\T J(q)^\T f.
\]
The production Jacobian is analytical; finite differences are used only for
verification.  Bending about both transverse axes, torsion, and a mixed field
are included.  All current plain, recovery-small, and recovery-finite cases
pass.  Their worst relative discrepancies are respectively
\[
  2.40\times10^{-9},\qquad
  2.24\times10^{-9},\qquad
  2.57\times10^{-9}.
\]

\section{FSI benchmark design}

The FSI runner supports NearestNeighbor, BeamMapper\_CoRotation,
BeamSplineMapper, and both recovery modes, together with scale, time step,
relaxation, coupling-iteration limit, kernel radius, regularization,
polynomial level, structural strategy, and target end time.  Each trial:
\begin{enumerate}
\item snapshots every modified JSON, MDPA, and material file;
\item synchronizes section and fluid-load scaling;
\item writes point history, CFD/CSM VTK, full log, settings, conclusion, and
      summary JSON into a unique directory;
\item monitors nonfinite output, coupling failure, and repeated structural
      nonconvergence;
\item stops invalid runs and verifies restoration by file hashes.
\end{enumerate}

The stability gate is \(0.5\rightarrow4.5\rightarrow10\rightarrow25\) s.
Failure of one gate prevents an unsupported longer run.  The analytical suite
remains the correctness reference.  In addition, the FSI runner uses a fixed,
physically scaled NearestNeighbor point-A \(u_x\) history as a trajectory
reference:
\[
 e_{\mathrm{NN}} =
 \sqrt{\frac{\sum_{t\in{\cal T}_{\mathrm{valid}}}
 (u_{\mathrm{mapper}}(t)-u_{\mathrm{NN}}(t))^2}
 {\sum_{t\in{\cal T}_{\mathrm{valid}}}u_{\mathrm{NN}}(t)^2}} .
\]
Here \({\cal T}_{\mathrm{valid}}\) is the exact intersection of finite sample
times, truncated at the mapper's last accepted time; no interpolation or
nonconverged output is used.  This metric compares nonlinear FSI trajectories,
but does not replace the prescribed-field analytical tests.

\section{Retained post-fix FSI results}

All retained trials use \(s=50\), \(\Delta t=0.05\), \(\alpha=0.03\), 60
coupling iterations, synchronized scaling, and \texttt{swap\_sign}.

\begin{center}
\begin{tabular}{p{40mm}p{15mm}p{18mm}p{22mm}p{43mm}}
\toprule
Mapper/mode & Target & Last valid & \(e_{\mathrm{NN}}\) & Failure\\
\midrule
NearestNeighbor reference
&4.5 s&4.50 s&0&completed\\
BeamSplineMapper, normalized
&4.5 s&3.05 s&\(1.0315\times10^{-2}\)&structural nonconvergence, step 62 at 3.10 s\\
Recovery, level 4, small
&4.5 s&4.20 s&\(3.8945\times10^{-2}\)&structural nonconvergence, step 85 at 4.25 s\\
Recovery, level 4, finite
&4.5 s&3.05 s&\(1.0242\times10^{-2}\)&structural nonconvergence, step 62 at 3.10 s\\
\bottomrule
\end{tabular}
\end{center}

The error column uses each case's own valid interval.  Since the small-recovery
case runs longer, a fair all-beam comparison also uses the shared
\(0.05\)--\(3.05\) s interval (61 samples): plain, recovery-small, and
recovery-finite respectively give
\[
 1.0315\times10^{-2},\qquad
 1.7555\times10^{-2},\qquad
 1.0242\times10^{-2}.
\]
Thus the small-mode value of \(3.8945\times10^{-2}\) over its longer
\(0.05\)--\(4.20\) s interval must not be directly ranked against the two
shorter values.

All four trial summaries report successful restoration.  The NN reference
completes 4.5 s, but no beam mapper reaches the 4.5 s gate, so no post-fix beam
run was promoted to 10 or 25 s.  The finite
recovery and normalized plain histories are almost identical at their limit:
their point-A \(u_x\) values at \(t=3.05\) s are \(0.01611986\) and
\(0.01612157\) m.  Recovery-small reproduces the earlier step-85 limit after
coordinate normalization and the Legendre polynomial correction.  Therefore
the old high-order monomial conditioning was not the root cause of that
event.

The analytical checks exclude coordinate mutation, a discontinuous projection
definition, missing Gaussian self-Hessian, a non-adjoint load map, and an
incorrect finite SO(3) tangent as current explanations.  FSI fields remain
finite up to the last accepted step, after which the structural nonlinear
solve repeatedly fails.  The evidence presently classifies the terminal event
primarily as a structural/coupled-trajectory nonconvergence, not as a mapper
singularity or mapper-produced NaN.  A structure-only replay of steps 62 and
85 is the next mechanism-isolating experiment.

\section{Reproducibility and retained records}

The canonical scripts are under \texttt{Scripts/}:
\begin{itemize}
\item \texttt{run\_fsi\_benchmark.py}: parameterized, restoring FSI runner;
\item \texttt{validate\_beam\_spline\_forward.py} and
      \texttt{validate\_beam\_spline\_adjoint.py};
\item \texttt{validate\_recovery\_forward.py} and
      \texttt{validate\_recovery\_adjoint.py};
\item the conditioning, Newton-history, VTK-history, and load-consistency
      analyzers.
\end{itemize}

Current analytical results are in
\texttt{TestCase\_Output/MapperVerification}.  Each validation script defines
its canonical output directory near the top of the file and, on rerun,
replaces only that exact direct child of \texttt{MapperVerification}.  This
keeps the normal invocation reproducible without accumulating versioned
copies.  Current nonlinear FSI evidence is restricted to the NN reference and
the three beam \(t=4.5\) gate directories under
\texttt{TestCase\_Output/StabilityTrials\_s50}.  The FSI runner avoids
aliasing NN point-A and point-B output streams and writes both own-valid-range
and common-range comparisons to
\texttt{nn\_rel\_l2\_comparison.json} and
\texttt{nn\_rel\_l2\_comparison.csv}.  Every numerical conclusion above can
be reconstructed from these summaries and the individual logs.

\section{Remaining theoretical decisions}

\begin{enumerate}
\item Confirm whether the desired recovery mapper is the paper-faithful,
      general three-dimensional Ahrem displacement field or a rigid-section
      finite-rotation beam surface model.  They are different nonlinear FSI
      models.
\item Confirm the precise semantics of the structural element's
      \texttt{ROTATION}: total exponential-coordinate vector, incremental
      rotation, or another composed parameter.  The finite hybrid requires a
      compatible total-state interpretation.
\item Confirm whether rigid cross-sections are a desired physical constraint.
      If section warping is relevant, the finite hybrid is too restrictive.
\item If a fully geometrically exact recovery is required, consider separate
      interpolation of translation and orientation on \(SO(3)\), with a
      consistent tangent transpose, rather than assigning finite physical
      meaning to \(\tfrac12\curl s\).
\end{enumerate}

\end{document}
